\documentclass[aspectratio=169,professionalfonts]{beamer}

\usetheme{Boadilla}
\usecolortheme{seahorse}
\setbeamertemplate{navigation symbols}{}
\setbeamertemplate{footline}[frame number]
\setbeamersize{text margin left=8mm,text margin right=8mm}

\usepackage[utf8]{inputenc}
\usepackage[T1]{fontenc}
\usepackage{lmodern}
\usepackage{amsmath,amssymb,bm,mathtools}
\usepackage{graphicx}
\usepackage{booktabs}
\usepackage{xcolor}
\usepackage{array}

\definecolor{PennBlue}{RGB}{1,31,91}
\definecolor{AccentRed}{RGB}{170,45,45}
\definecolor{SoftGray}{RGB}{245,245,245}
\setbeamercolor{title}{fg=PennBlue}
\setbeamercolor{frametitle}{fg=PennBlue}
\setbeamercolor{block title}{fg=white,bg=PennBlue}
\setbeamercolor{block body}{bg=SoftGray}
\setbeamercolor{alerted text}{fg=AccentRed}

\newcommand{\E}{\mathbb{E}}
\newcommand{\Prb}{\mathbb{P}}
\newcommand{\Var}{\operatorname{Var}}
\newcommand{\Cov}{\operatorname{Cov}}
\newcommand{\Corr}{\operatorname{Corr}}
\newcommand{\sd}{\operatorname{sd}}
\newcommand{\logit}{\operatorname{logit}}
\newcommand{\logistic}{\operatorname{logistic}}
\newcommand{\cost}{\operatorname{cost}}
\newcommand{\Normal}{\mathcal{N}}
\newcommand{\Bern}{\operatorname{Bernoulli}}
\newcommand{\BetaDist}{\operatorname{Beta}}
\newcommand{\Unif}{\operatorname{Unif}}
\newcommand{\QB}{\mathrm{QB}}
\newcommand{\notQB}{\mathrm{not\text{-}QB}}

% The Loser's Curse figures referenced in body.tex were not uploaded.
% This macro lets the deck compile now and automatically uses the figure
% when the corresponding file is placed in the expected path.
\newcommand{\maybefig}[2][]{%
  \IfFileExists{#2}{\includegraphics[#1]{#2}}{%
    \fbox{\begin{minipage}[c][0.42\textheight][c]{0.90\linewidth}
      \centering\scriptsize Missing figure file:\\[1mm]
      \texttt{\detokenize{#2}}\\[2mm]
      Put this file in the same relative path to render the final deck.
    \end{minipage}}%
  }%
}

\newcommand{\slidecite}[1]{\vspace{0.5mm}{\scriptsize\hfill #1}}

\title{Three Problems in Sports Analytics}
\author{Professor Abraham (Adi) Wyner  }
\institute{Tel Aviv University Statistics Department}
\date{June 16, 2026}

\begin{document}

\begin{frame}
\frametitle{Three Problems in Sports Analytics}
\begin{columns}[T,totalwidth=\textwidth]
\begin{column}{0.52\textwidth}
  {\Large Professor Abraham (Adi) Wyner}\\[1mm]
  Tel Aviv University Statistics Department\\
  June 16, 2026
  \vspace{2mm}
  
  {\small Joint work with students\\[0.5mm]
  Ryan Brill \& J Pipping}


  \vspace{4mm}
  \begin{block}{Sports analytics}
  The use of data, statistics, and mathematical modeling to understand games, evaluate people, forecast outcomes, and make better decisions.
  \end{block}
\end{column}
\begin{column}{0.44\textwidth}
  \begin{block}{Four pillars}
  \begin{enumerate}
    \item Historical descriptions and evaluations of teams and players.
    \item Prediction: future performance and evaluation.
    \item Strategic decisions.
    \item Training analytics.
  \end{enumerate}
  \end{block}
\end{column}
\end{columns}
\end{frame}

\begin{frame}
\frametitle{Why sports analytics is booming}
\begin{columns}[T,totalwidth=\textwidth]
\begin{column}{0.49\textwidth}
\begin{block}{The demand side}
\begin{itemize}
  \item Roughly a \$40 billion industry.
  \item Fans, owners, players, media, and gamblers care deeply about the games.
  \item The questions are simple to state and hard to answer well.
\end{itemize}
\end{block}
\end{column}
\begin{column}{0.49\textwidth}
\begin{block}{The intellectual side}
\begin{itemize}
  \item A laboratory for behavior: decisions, incentives, psychology, health, and race.
  \item Hard statistics: selection, censoring, calibration, causality, and high-dimensional prediction.
  \item Hard math: stochastic processes, optimization, networks, game theory, and tails.
\end{itemize}
\end{block}
\end{column}
\end{columns}

\vspace{2mm}
\centering

\alert{This talk has something for all four: sports, human behavior, statistics, and mathematics.}
\end{frame}

\begin{frame}
\frametitle{The three problems}
\begin{enumerate}
  \item \textbf{Hot Hand Fallacy.} We see streaks and infer momentum. Are hot streaks real, or can finite samples fool us?
  \item \textbf{Paradox of the Blown Lead.} We think teams blow big leads often. Are win probability models wrong, or are we conditioning on the wrong event?
  \item \textbf{Loser's Curse in the NFL Draft.} Big-time decision makers trade valuable assets for draft picks. Are they making poor decisions, or optimizing a different objective?
\end{enumerate}

\end{frame}

\begin{frame}
\frametitle{A common theme: Signal is hard to spot}


\large
\begin{enumerate}
\item It is not so easy to estimate ``hotness.''
\item Blown leads are more common than you think. \\[1mm]
\item Are high-value decisions just as susceptible to psychological biases as small-value decisions?
\end{enumerate}


\vspace{5mm}
\begin{block}{Statistical moral}
Before calling behavior biased or irrational, specify the estimand and the utility function.
\end{block}
\end{frame}

\begin{frame}
\frametitle{Problem I: The Loser's Curse in the NFL Draft}
\begin{block}{Why this example matters}
The NFL draft is not a toy experiment. It is a repeated, high-stakes labor market in which teams reveal their valuations by trading picks.
\end{block}

\begin{columns}[T,totalwidth=\textwidth]
\begin{column}{0.49\textwidth}
\begin{itemize}
  \item Worse teams pick earlier.
  \item Teams trade picks for picks, future picks, and players.
  \item Draft order has prices, costs, uncertain payoffs, and strategic value.
\end{itemize}
\end{column}
\begin{column}{0.49\textwidth}
\begin{block}{Behavioral-economics test}
Massey--Thaler ask whether expert decision makers still make biased decisions when the \emph{``monetary stakes are exceedingly high.''}
\slidecite{Massey and Thaler (2013)}
\end{block}
\end{column}
\end{columns}

\vspace{2mm}
\centering
\alert{If NFL general managers fail here, the behavioral story is much bigger than a lab curiosity.}
\end{frame}

\begin{frame}
\frametitle{Massey--Thaler's Two Claims}
\begin{enumerate}
  \item \textbf{The Loser's Curse.} The NFL rewards last year's losers with the top draft picks, but those picks are not the best bargains under expected surplus value.
  \item \textbf{The trade market is too steep.} The trade valuation curve looks much like the old Jimmy Johnson chart, while realized production value declines much more slowly with draft position.
\end{enumerate}

\vspace{3mm}
\begin{block}{The provocative conclusion}
If the right utility function is expected on-field production per salary dollar, then teams that trade up are burning value and teams should trade down more often.
\end{block}
\end{frame}

\begin{frame}
\frametitle{Claim 1: The Gift that Becomes a Curse}
Let $x$ denote draft position, $Y$ first-contract on-field production value, and $c(x)$ rookie-contract cost.
\[
  \text{expected surplus value} \;=\; \E\{Y-c(x)\mid x\}.
\]

\begin{columns}[T,totalwidth=\textwidth]
\begin{column}{0.52\textwidth}
\begin{itemize}
  \item The worst teams get the smallest $x$: the right to choose first.
  \item Top picks have the highest expected raw production.
  \item But they also have large costs and large forecasting error.
  \item Massey--Thaler's striking result: surplus per dollar is highest near the \alert{end of the first round}, not the top.
\end{itemize}
\end{column}
\begin{column}{0.44\textwidth}
\centering
\maybefig[width=\linewidth,height=0.48\textheight,keepaspectratio]{plots/plot_MasseyThaler_finding1.png}
\end{column}
\end{columns}

\vspace{1mm}
\centering
\Large \alert{Ouch: the league's reward may be a poisoned gift.}
\end{frame}

\begin{frame}
\frametitle{What Makes it a ``Curse''?}
\begin{block}{The allocation rule}
The worst teams are placed at the front of the line, precisely where expected surplus value is low relative to the opportunity cost.
\end{block}

\vspace{1mm}
\begin{itemize}
  \item The policy goal is competitive balance: give bad teams access to the best amateurs.
  \item The economic question is different: is pick 1 the best use of scarce salary-cap resources?
  \item Under expected surplus, Massey--Thaler's answer is no.
\end{itemize}

\vspace{2mm}
\begin{block}{Their sharp version}
\centering
\Large
``the first pick in the first round is the worst pick in the round.''
\slidecite{Massey and Thaler (2013)}
\end{block}
\end{frame}

\begin{frame}
\frametitle{Claim 2: The Trade Market is Too Steep}
\begin{columns}[T,totalwidth=\textwidth]
\begin{column}{0.48\textwidth}
\begin{block}{Market value curve}
Estimate $M(x)$ from observed pick-for-pick trades:
\[
  \sum_{x\in A} M(x) \approx \sum_{x\in B} M(x).
\]
Empirically, this curve closely mirrors the traditional Jimmy Johnson trade chart.
\end{block}
\end{column}
\begin{column}{0.48\textwidth}
\begin{block}{Production/surplus curve}
Estimate $V(x)$ from realized player production and compensation:
\[
  V(x) = \E\{Y-c(x)\mid x\}.
\]
This curve declines much more slowly than the market curve.
\end{block}
\end{column}
\end{columns}

\vspace{2mm}
\centering
\alert{Same pick. Two value functions. Very different prices.}
\end{frame}

\begin{frame}
\frametitle{The Picture: Market Value vs. Surplus Value}
\begin{columns}[T,totalwidth=\textwidth]
\begin{column}{0.49\textwidth}
\centering
\maybefig[width=\linewidth,height=0.47\textheight,keepaspectratio]{plots/plot_MasseyThaler_finding2.png}
\end{column}
\begin{column}{0.49\textwidth}
\begin{itemize}
  \item The market pays a huge premium for moving up.
  \item Expected surplus does not fall nearly that fast.
  \item Trading down converts one glamorous pick into several cheaper expected-surplus opportunities.
\end{itemize}

\vspace{2mm}
\begin{block}{Massey--Thaler bottom line}
\centering
``top draft picks are significantly overvalued.''
\slidecite{Massey and Thaler (2013)}
\end{block}
\end{column}
\end{columns}
\end{frame}

\begin{frame}
\frametitle{A One-Line Trade-Up Diagnosis}
Suppose a team trades a bundle $B$ for an earlier pick $a$.

\vspace{1mm}
\begin{align*}
\text{Market says fair:} \qquad & M(a) \approx \sum_{x\in B} M(x),\\[1mm]
\text{Expected surplus says:} \qquad & \E\{S\mid a\} \ll \sum_{x\in B}\E\{S\mid x\}.
\end{align*}

\vspace{2mm}
\begin{block}{If expected surplus is the utility function}
The team trading up is paying market price for a dream and giving away the better expected bargain.
\end{block}

\vspace{2mm}
\centering
\Large
\alert{The obvious prescription: trade down. A lot.}
\end{frame}

\begin{frame}
\frametitle{The Behavioral-Economics Punchline}
\begin{columns}[T,totalwidth=\textwidth]
\begin{column}{0.50\textwidth}
\begin{block}{Why economists cared}
This is a rare high-value setting with experienced experts, repeated feedback, public prices, and consequences for jobs, profits, and championships.
\end{block}
\end{column}
\begin{column}{0.46\textwidth}
\begin{block}{The interpretation}
If teams keep trading up despite the surplus evidence, then perhaps laboratory-style biases scale up to the NFL front office.
\end{block}
\end{column}
\end{columns}

\vspace{3mm}
\begin{center}
\Large
\alert{Are million-dollar decisions just as psychologically fragile as small-stakes lab decisions?}
\end{center}
\end{frame}

\begin{frame}
\frametitle{Twelve-Plus Years Later: Did Teams Learn?}
\begin{block}{The natural follow-up}
Massey asked Ryan Brill to update the analysis with modern data through 2025.
\end{block}

\begin{columns}[T,totalwidth=\textwidth]
\begin{column}{0.52\textwidth}
\begin{itemize}
  \item The rookie wage scale changed.
  \item Analytics departments grew.
  \item The Massey--Thaler result became famous.
  \item Teams had many drafts to learn from their own outcomes.
\end{itemize}
\end{column}
\begin{column}{0.44\textwidth}
\begin{block}{The awkward answer}
The broad picture does \alert{not} go away: the trade market remains steep and the expected-surplus critique persists.
\end{block}
\end{column}
\end{columns}

\vspace{2mm}
\centering
\Large
\alert{Provocative reading: ``just as stupid as ever.''}
\end{frame}

\begin{frame}
\frametitle{The Replication Sting}
\centering
\maybefig[width=0.88\linewidth,height=0.56\textheight,keepaspectratio]{plots/plot_Massey_Thaler_replicate.png}

\vspace{1mm}
\begin{block}{Why this is such a strong behavioral story}
If expected surplus is the correct objective, then the NFL draft market is not merely biased; it is biased in a way that survived publication, money, feedback, and a decade of analytics.
\end{block}
\end{frame}

\begin{frame}
\frametitle{But Maybe Football People Are Not So Stupid}
\begin{center}
\Large
What if the mistake is not in the front office, but in the statistician's utility function?
\end{center}

\vspace{2mm}
\begin{block}{Brill--Wyner move}
Do not assume teams maximize expected surplus. Ask whether their trades are consistent with a different objective: acquiring \alert{top talent}.
\end{block}

\begin{itemize}
  \item Championships are nonlinear outcomes.
  \item A roster is not built from expected value alone.
  \item One elite quarterback can matter more than several merely good bargains.
  \item The draft may be a search for the right tail, not the mean.
\end{itemize}
\end{frame}

\begin{frame}
\frametitle{The Decision Criterion Changes}
Traditional expected-surplus analysis asks for
\[
  \E\{Y-c(x)\mid x\}.
\]
The top-talent hypothesis asks for tail probabilities such as
\[
  \Prb\{Y>r\mid x\}
  \quad\text{or}\quad
  \Prb\{Y-c(x)>r\mid x\}.
\]

\vspace{2mm}
\begin{block}{Statistical moral}
Before calling a decision biased, specify the utility function. A mean regression answers a mean question; it does not answer a championship question.
\end{block}
\end{frame}

\begin{frame}
\frametitle{A Concrete Motivating Trade}
\begin{block}{2013: Miami moves up to pick 3}
Miami gave picks 12 and 42 to Oakland for pick 3.
\end{block}

\begin{center}
\begin{tabular}{lcc}
\toprule
Valuation rule & Winner under the rule & Magnitude\\
\midrule
Jimmy Johnson chart & Miami & Miami draft capital $+31\%$\\
Expected surplus value & Oakland & Oakland draft capital $+135\%$\\
Expected elite-player count & Miami & Miami draft capital $+40\%$\\
\bottomrule
\end{tabular}
\end{center}

\vspace{2mm}
\small
Same trade. Same draft board. Different objective functions imply different winners.
\end{frame}

\begin{frame}
\frametitle{Why Expectation May Be the Wrong Utility}
A general manager is not literally trying to maximize the average second-contract value of drafted players. The strategic objective may be closer to
\[
  \Prb\{\text{build a Super Bowl-winning roster}\}.
\]

\begin{block}{One-line threshold example}
Let team performance be random. If
\[
A\sim \Normal(1,0),\qquad B\sim \Normal(0,1),\qquad \tau=1.28,
\]
then
\[
\Prb(A>\tau)=0,\qquad \Prb(B>\tau)\approx 0.10.
\]
The lower-mean, higher-variance option can be better for exceeding a high threshold.
\end{block}
\end{frame}

\begin{frame}
\frametitle{Draft Picks as Right-Tail Lotteries}
\begin{block}{Hypothesis}
Teams trade for variance and right-tail outcomes, not only for expected value.
\end{block}

\begin{itemize}
  \item A good roster needs many competent players.
  \item A championship roster may require a few transformational players.
  \item Top draft picks may be valuable because they are more likely to produce extreme right-tail outcomes.
\end{itemize}

\vspace{2mm}
\begin{center}
\Large
Estimate the \alert{full conditional distribution} $Y\mid x$, not only $\E[Y\mid x]$.
\end{center}
\end{frame}

\begin{frame}
\frametitle{Empirical Clue: Both Mean and Variance Decline}
\centering
\maybefig[width=0.92\linewidth,height=0.60\textheight,keepaspectratio]{plots/plot_empMeanSd.png}

\vspace{1mm}
\small
If high picks have both higher mean and higher variance, expected value can understate their value for threshold objectives.
\end{frame}

\begin{frame}
\frametitle{Measuring Player Performance}
\begin{block}{Outcome variable}
$Y$ is a player's second-contract average salary per year as a percentage of the salary cap.
\end{block}

\begin{itemize}
  \item Interprets the free-agent market as an approximate valuation mechanism.
  \item If a player does not sign a second contract, set $Y=0$.
  \item Focuses on first-contract performance and first-contract surplus.
\end{itemize}

\begin{block}{Caveat}
This is not a perfect player value measure. The point is to compare draft-pick value functions given a performance measure.
\end{block}
\end{frame}

\begin{frame}
\frametitle{A Bayesian Spike-Plus-Beta Model}
Let $y_{\rm bust}=0.5\%$ of the salary cap. Model a spike near zero and a continuous right tail.

\vspace{2mm}
\begin{align*}
1\{Y\le y_{\rm bust}\}\mid x &\sim \Bern\{p_b(x)\},
& p_b(x)&=\logistic(\alpha_0+\alpha_1x),\\
Y\mid x,Y>y_{\rm bust} &\sim \BetaDist\{\mu(x),\phi(x)\},
& \mu(x)&=\logistic(\widetilde x^\top\beta),\\
&& \phi(x)&=\exp(\gamma_0+\gamma_1x).
\end{align*}

\vspace{1mm}
Weakly informative priors; posterior inference by Stan MCMC.
\end{frame}

\begin{frame}
\frametitle{What the Density Model Learns}
\begin{columns}[T,totalwidth=\textwidth]
\begin{column}{0.54\textwidth}
\begin{itemize}
  \item The bust probability increases with draft position.
  \item The right tail shifts left as the draft progresses.
  \item The distribution changes shape; it is not just a moving mean.
\end{itemize}
\end{column}
\begin{column}{0.44\textwidth}
\centering
\maybefig[width=\linewidth,height=0.50\textheight,keepaspectratio]{plots/plot_post_density_full_rdsall_SE.png}
\end{column}
\end{columns}

\vspace{2mm}
\begin{block}{Why this matters}
Any utility based on elite outcomes needs $\Prb(Y>r\mid x)$, not just $\E[Y\mid x]$.
\end{block}
\end{frame}

\begin{frame}
\frametitle{Position-Specific Hierarchy}
The model is extended to condition on player position $p$.

\vspace{1mm}
\begin{align*}
Y\mid x,p,Y>y_{\rm bust} &\sim \BetaDist\{\mu(x,p),\phi(x,p)\},\\
\mu(x,p)&=\logistic(\widetilde x^\top\beta_p),\\
\phi(x,p)&=\exp(\gamma_{0,p}+\gamma_{1,p}x),\\
p_b(x,p)&=\logistic(\alpha_{0,p}+\alpha_{1,p}x).
\end{align*}

\begin{block}{Why hierarchical?}
There are many position groups and sparse data at each position-by-pick combination, so position-specific parameters are shrunk toward overall means.
\end{block}
\end{frame}

\begin{frame}
\frametitle{Right-Tail Draft Value Curves}
For an elite-performance threshold $r$, define
\[
  v_r(x)=\frac{\Prb(Y>r\mid x)}{\Prb(Y>r\mid x=1)}.
\]

\begin{block}{Interpretation}
$v_r(x)$ is proportional to the expected number of elite players produced by pick $x$, relative to pick 1.
\end{block}

\begin{center}
\begin{tabular}{ll}
\toprule
Threshold $r$ & Reference level in the paper\\
\midrule
$24.5\%$ & Joe Burrow level\\
$19.7\%$ & Deshaun Watson level\\
$17.8\%$ & Jared Goff level\\
$15.0\%$ & Derek Carr level\\
$11.0\%$ & Penei Sewell level\\
$8.7\%$ & Tre'Davious White level\\
\bottomrule
\end{tabular}
\end{center}
\end{frame}

\begin{frame}
\frametitle{Main Result: The Tail Curves are Steep}
\centering
\maybefig[width=0.92\linewidth,height=0.60\textheight,keepaspectratio]{plots/plot_tail_probs_relative_SE.png}

\vspace{1mm}
\small
As the elite threshold $r$ increases, $x\mapsto \Prb(Y>r\mid x)$ steepens and can approach the fitted trade market curve.
\end{frame}

\begin{frame}
\frametitle{Revealed Preference Interpretation}
\begin{block}{Result}
The fitted trade market is closest, in mean absolute error, to a right-tail curve with threshold
\[
  r\approx 19.7\% \text{ of the salary cap}.
\]
\end{block}

\begin{itemize}
  \item This is roughly Deshaun Watson's second-contract level in the paper's scale.
  \item It is a very high bar: among regular-starter quarterbacks in the dataset, about $8/65=12.3\%$ met it.
  \item This does not prove what general managers think; it says their trades are consistent with valuing extreme upside.
\end{itemize}
\end{frame}

\begin{frame}
\frametitle{Surplus Value, But in the Right Tail}
Let $S=Y-c(x)$. Massey--Thaler focus on $\E[S\mid x]$. The alternative is
\[
  \Prb(S>r\mid x).
\]

\centering
\maybefig[width=0.92\linewidth,height=0.50\textheight,keepaspectratio]{plots/plot_G_surplusValueCurves_2A.png}

\vspace{1mm}
\small
The expected-surplus first-round spike is much less decisive when the objective is a high-surplus tail event.
\end{frame}

\begin{frame}
\frametitle{Quarterbacks are Different}
\begin{columns}[T,totalwidth=\textwidth]
\begin{column}{0.46\textwidth}
\begin{itemize}
  \item Quarterbacks have higher expected value.
  \item More importantly, quarterbacks have much fatter right tails.
  \item Top-pick quarterback upside is the central economic force in the draft market.
\end{itemize}
\end{column}
\begin{column}{0.52\textwidth}
\centering
\maybefig[width=\linewidth,height=0.55\textheight,keepaspectratio]{plots/plot_G_valueCurves_byQB.png}
\end{column}
\end{columns}
\end{frame}

\begin{frame}
\frametitle{For Elite QBs, the Curse Largely Disappears}
\centering
\maybefig[width=0.92\linewidth,height=0.58\textheight,keepaspectratio]{plots/plot_G_surplusValueCurves_byQB.png}

\vspace{1mm}
\small
For elite quarterback surplus, the top of the draft can be rationally valuable despite its high cost. For non-quarterbacks, the first-round surplus peak persists.
\end{frame}

\begin{frame}
\frametitle{Revisiting the 2013 Dolphins--Raiders Trade}
\begin{center}
\begin{tabular}{lccc}
\toprule
Objective & Pick 3 & Picks 12+42 & Preferred side\\
\midrule
Johnson chart & 2200 & 1680 & Pick 3\\
Expected surplus & 1.02 & 2.40 & Picks 12+42\\
Expected elite players & 0.893 & 0.639 & Pick 3\\
\bottomrule
\end{tabular}
\end{center}

\begin{block}{Lesson}
``Overpaying'' is not an intrinsic property of a trade. It is relative to a utility function.
\end{block}
\end{frame}

\begin{frame}
\frametitle{What Survives of the Loser's Curse?}
\begin{columns}[T,totalwidth=\textwidth]
\begin{column}{0.49\textwidth}
\begin{block}{Under expected surplus}
Top picks are expensive, and the best expected bargains are later in the first round.
\end{block}
\end{column}
\begin{column}{0.49\textwidth}
\begin{block}{Under elite-player probability}
Top picks can be correctly priced because they buy access to the far right tail.
\end{block}
\end{column}
\end{columns}

\vspace{4mm}
\centering
\Large
The empirical dispute is not only about estimates.\newline
It is about the decision criterion.
\end{frame}

\begin{frame}
\frametitle{Methodological Takeaway}
A mean regression gives one value curve. A conditional density gives many.

\vspace{2mm}
\begin{align*}
\E[Y\mid x] &\quad \text{performance value},\\
\E[Y-c(x)\mid x] &\quad \text{surplus value},\\
\Prb(Y>r\mid x) &\quad \text{elite-player value},\\
\Prb(Y-c(x)>r\mid x) &\quad \text{elite-surplus value}.
\end{align*}

\begin{block}{General principle}
Estimate the distribution first; choose the functional after the decision problem is specified.
\end{block}
\end{frame}

\begin{frame}
\frametitle{Strategic Takeaway for Teams}
\begin{itemize}
  \item The paper is not a proof that teams should always trade up.
  \item It is evidence that top picks may be rationally expensive for teams seeking transformational upside.
  \item The optimal action depends on team state: roster, current quarterback, cap, injuries, prospect pool, and future picks.
\end{itemize}

\vspace{2mm}
\begin{block}{Ultimate target}
\[
  \arg\max_{\text{trade/action }a}\; \E\{U(\text{championship path},\text{ roster},\text{ cap})\mid a\}.
\]
\end{block}
\end{frame}



%\begin{frame}
%\frametitle{Problem II: The Hot Hand Fallacy}
%\begin{block}{Empirical question}
%Do players shoot better after a streak of made shots?
%\end{block}

%\vspace{1mm}
%A natural statistic is a conditional success rate:
%\[
 % \widehat P(S\mid S^k)
 % =
%  \frac{\#\{\text{successes following }k\text{ successes}\}}
 %      {\#\{\text{occurrences of }k\text{ successes}\}}.
%\]

%\vspace{1mm}
%For basketball this feels natural because shots are ordered, repeated, and naturally split into games.

%\vspace{2mm}
%\begin{block}{The trap}
%The statistic is simple; the averaging problem is not.
%\end{block}
%\end{frame}


\begin{frame}{Hot Hand Fallacy: Motivating Example}
\large

A basketball player is known to be a  \(50\%\) free throw shooter.  The player wants to know whether he gets ``hot'' by making three shots in a row.  For each game, he computes
\[
r_i
=
\frac{\#\{\text{makes following three consecutive makes in game }i\}}
{\#\{\text{three-consecutive-make opportunities in game }i\}}.
\]

\medskip

He then averages these game-level ratios across games and finds
\[
\frac{1}{G}\sum_{i=1}^{G} r_i = 0.46.
\]

\medskip

\begin{block}{Question}
This appears to be evidence against hotness. Is it?
\end{block}

\end{frame}

\begin{frame}{Second Motivating Example}
\large

Now consider a different version of the same question.
The player takes \(100\) three-point shots over the season and makes exactly
\(50\) of them.  

\medskip

Here the natural null model is not iid Bernoulli with known \(p\), but the
fixed-\(m\) permutation model: the observed sequence is treated as a uniform
random permutation of
\[
50 \text{ makes}
\qquad\text{and}\qquad
50 \text{ misses}.
\]

\medskip

He looks across the full sequence of \(100\) shots and finds that his shooting
percentage after three consecutive makes is 
\[
0.44.
\]

\medskip

\begin{block}{Question}
What should he make of this? Evidence for or against hotness?
\end{block}

\end{frame}

\begin{frame}{The hot-hand question}
\begin{block}{Empirical question}
Do players shoot better after a streak of made shots?
\end{block}

\vspace{0.2cm}
A natural statistic is a conditional success rate:
\[
  \widehat P(S\mid S^k)
  =
  \frac{\#\{\text{successes following }k\text{ successes}\}}
       {\#\{\text{occurrences of }k\text{ successes}\}}.
\]

\vspace{0.2cm}
For basketball, this feels natural because shots are:
\begin{itemize}
  \item ordered,
  \item homogeneous enough to model simply, especially free throws,
  \item and naturally split into games.
\end{itemize}


\end{frame}

\begin{frame}{Motivation: Gilovich--Vallone--Tversky}
\begin{itemize}
  \item The hot-hand question is simple and compelling:
  \[
    \text{Do players shoot better after a streak of made shots?}
  \]
  \item Gilovich, Vallone, and Tversky (1985) argued that the hot hand was largely a
  \textbf{misperception of random sequences}.
  \item Their analysis relied heavily on conditional success rates after makes and misses:
  \[
    \widehat P(S \mid S^k)
    =
    \frac{\#\{\text{successes following } k \text{ successes}\}}
         {\#\{\text{occurrences of } k \text{ successes}\}}.
  \]
  \item In their data, average success rates after streaks of makes appeared close to the player's
  base success rate, suggesting little evidence of positive serial dependence.
\end{itemize}

\vspace{0.2cm}
\begin{block}{The original interpretation}
Fans see streaks; the conditional-percentage evidence seemed to say:
\[
  \text{what looks like a hot hand may just be randomness.}
\]
\end{block}
\end{frame}

\begin{frame}{Miller--Sanjurjo: The Rebuttal}
\begin{itemize}
  \item Miller and Sanjurjo (2018) identified a finite-sample problem with the standard
  conditional-percentage calculation.
  \item Their \textbf{hot-hand paradox}: even for iid coin tosses,
  the expected proportion of heads following a string of heads in a finite sequence is
  \textbf{less than} \(1/2\).
  \item Thus, observing
  \[
    \widehat P(S \mid S^k) \approx \widehat P(S)
  \]
  need not mean ``no hot hand.'' The statistic can be biased downward under the no-hot-hand null.
\end{itemize}

\vspace{0.1cm}
\begin{block}{Miller--Sanjurjo insight}
Even when shots are iid, the naive conditional percentage after a streak can be biased downward.
\end{block}


\end{frame}

\begin{frame}{A Covariance View of the Hot-Hand Bias}
\begin{itemize}
  \item This talk does \emph{not} reprove the hot-hand literature.
  \item The goal is to give a simple structural explanation for why the naive statistic can be biased.
  \item The explanation is not psychological; it is algebraic:
\end{itemize}

\vspace{0.35cm}
\[
  \boxed{
  \text{bias}
  =
  \text{denominator--rate covariance}
  }
\]

\vspace{0.35cm}
\begin{block}{Guiding idea}
When we average game-level conditional percentages, games with few qualifying streaks
can receive the same weight as games with many qualifying streaks. The resulting mean of ratios
need not equal the pooled conditional probability.
\end{block}
\end{frame}

\begin{frame}{The basic setup}
Let
\[
X_1,\dots,X_L \overset{\text{iid}}{\sim} \text{Bernoulli}(p),
\qquad X_t=1 \text{ means success}.
\]
Fix a run length $k$. Define
\[
D=d_{L,k}^{+}
=
\#\{\text{$k$-success runs starting in positions }1,\dots,L-k\},
\]
\[
N=n_{L,k}^{+}
=
\#\{\text{$k$-success runs followed by another success}\},
\]
and the within-sequence conditional rate
\[
R=r_{L,k}^{+}=\frac{N}{D}\quad(D>0).
\]

\vspace{0.2cm}
Under the iid null, the true target is
\[
P(S\mid S^k)=p.
\]
\end{frame}

\begin{frame}{Two ways to average across games}
Suppose we observe $G$ independent games or sequences.

\vspace{0.2cm}
\begin{columns}[T]
\begin{column}{0.48\textwidth}
\begin{block}{Equal-weight average}
\[
\bar R_G=\frac{1}{G}\sum_{i=1}^G R_i.
\]
Each game gets the same weight.
\end{block}
\end{column}
\begin{column}{0.48\textwidth}
\begin{block}{Pooled estimator}
\[
\widetilde R_G=
\frac{\sum_{i=1}^G N_i}{\sum_{i=1}^G D_i}
=\sum_{i=1}^G
\frac{D_i}{\sum_j D_j}R_i.
\]
Each opportunity gets the same weight.
\end{block}
\end{column}
\end{columns}

\vspace{0.3cm}
\begin{alertblock}{Key distinction}
A game with one qualifying streak should not carry the same information as a game with twenty qualifying streaks.
\end{alertblock}
\end{frame}

\begin{frame}{Main identity: the whole issue in one line}
Since $DR=N$,
\[
  \E[N]=\E[DR]=\E[D]\E[R]+\Cov(R,D).
\]
Under iid Bernoulli sampling,
\[
  \E[N]=p\E[D].
\]
Therefore
\[
\boxed{
  \E[R]
  =
  p-\frac{\Cov(R,D)}{\E[D]}
}
\]

\vspace{0.3cm}
\begin{block}{Interpretation}
If $R$ and $D$ are positively correlated, then
\[
  \E[R]<p.
\]
The equal-weight average of game-level ratios is biased downward.
\end{block}
\end{frame}

\begin{frame}{The heuristic: why should $R$ and $D$ be positively correlated?}
For the $S^k$ pattern, a success after the run does two things:
\[
  S^k \mapsto S^{k+1}.
\]

\vspace{0.2cm}
It contributes to the numerator $N$ and extends the run, creating more overlapping $k$-success opportunities.

\vspace{0.3cm}
For example, with $k=3$:
\[
  SSSS S
\]
contains multiple overlapping $SSS$ opportunities.

\vspace{0.3cm}
\begin{alertblock}{Core intuition}
Sequences with many opportunities $D$ tend to be sequences where successes continue after success runs. Those same sequences tend to have larger $R=N/D$.
\end{alertblock}
\end{frame}

\begin{frame}{Consequences under iid Bernoulli sampling}
For fixed $k$ and $p\in(0,1)$, the first-order bias after $k$ successes is
\[
  \E[r_{L,k}^{+}]
  =
  p-\frac{p^{1-k}-p}{L-k}+o(L^{-1}).
\]
Thus success-after-success-runs is biased downward.

\vspace{0.2cm}
Similarly, after $k$ failures,
\[
  \E[r_{L,k}^{-}]
  =
  p+\frac{(1-p)^{1-k}-(1-p)}{L-k}+o(L^{-1}).
\]
Thus success-after-failure-runs is biased upward.

\vspace{0.2cm}
So the naive contrast has negative mean even when the true contrast is zero:
\[
  \E[r_{L,k}^{+}-r_{L,k}^{-}]
  =
  -\frac{p^{1-k}+(1-p)^{1-k}-1}{L-k}+o(L^{-1}).
\]
\end{frame}

\begin{frame}{Simulation: iid Bernoulli, $k=3$, $p=1/2$}
\centering
\maybefig[width=0.98\textwidth,height=0.67\textheight,keepaspectratio]{plots/two_panel.png}
\end{frame}

\begin{frame}{What the iid simulation shows}
For $k=3$, $p=1/2$, and 10,000 sequences:

\vspace{0.2cm}
\begin{columns}[T]
\begin{column}{0.48\textwidth}
\begin{block}{$L=30$}
\[
\begin{aligned}
\widehat{\Cov}(r,d)&\approx 0.444,\\
\widehat{\Corr}(r,d)&\approx 0.668,\\
\bar r&\approx 0.384,\\
\bar r+\widehat{\Cov}(r,d)/\bar d&\approx 0.503.
\end{aligned}
\]
\end{block}
\end{column}
\begin{column}{0.48\textwidth}
\begin{block}{$L=100$}
\[
\begin{aligned}
\widehat{\Cov}(r,d)&\approx 0.480,\\
\widehat{\Corr}(r,d)&\approx 0.630,\\
\bar r&\approx 0.457,\\
\bar r+\widehat{\Cov}(r,d)/\bar d&\approx 0.497.
\end{aligned}
\]
\end{block}
\end{column}
\end{columns}

\vspace{0.2cm}
\begin{alertblock}{Message}
The pooled estimator recovers $p\approx 1/2$; the equal-weight game average is lower because $r$ and $d$ are positively correlated.
\end{alertblock}
\end{frame}

\begin{frame}{Basketball interpretation}
\small
A player looking back over a season might compute
\[
  \text{game-by-game } P(S\mid S^k),
  \quad \text{then average over games.}
\]
That feels natural because games are the natural record-keeping units.

\vspace{0.15cm}
But the target is the rate over \emph{opportunities}, not over games:
\[
  \frac{\sum_i N_i}{\sum_i D_i}.
\]
This is what the player would get by ignoring artificial game endings and stringing all games into one long season sequence.

\vspace{0.15cm}
\begin{block}{Gap between the two methods}
\[
  p-\E[R]
  =
  \frac{\Cov(R,D)}{\E[D]}.
\]
Positive covariance means the game-level average is systematically too low.
\end{block}
\end{frame}

\begin{frame}{Extension: fixed-$m$ permutation model}
Now suppose the number of successes is fixed.

\vspace{0.2cm}
Among $L$ trials, exactly $m$ are successes, and all
\[
  \binom{L}{m}
\]
sequences are equally likely. Let
\[
  p=\frac{m}{L}.
\]

\vspace{0.25cm}
This is the natural null if we condition on a player's observed number of makes.

\vspace{0.3cm}
The pooled finite-population target is not exactly $p$:
\[
  \frac{\E[N]}{\E[D]}
  =
  \frac{m-k}{L-k}
  =
  p-\frac{k(1-p)}{L-k}.
\]

\vspace{0.15cm}
So there is now a without-replacement correction even before the ratio issue.
\end{frame}

\begin{frame}{Permutation decomposition}
In the fixed-$m$ permutation model,
\[
\boxed{
  p-\E[r_{L,k}]
  =
  \frac{k(1-p)}{L-k}
  +
  \frac{\Cov(r_{L,k},d_{L,k})}{\E[d_{L,k}]}
}
\]

\vspace{0.2cm}
This separates the bias into two pieces:
\begin{enumerate}
  \item \textbf{Finite-population term:}
  \[
    \frac{k(1-p)}{L-k}.
  \]
  This comes from sampling without replacement.

  \item \textbf{Covariance term:}
  \[
    \frac{\Cov(r_{L,k},d_{L,k})}{\E[d_{L,k}]}.
  \]
  This is the same denominator--rate mechanism as in the Bernoulli model.
\end{enumerate}

\vspace{0.1cm}
For $k=1$, the covariance term vanishes exactly. For $k\ge 2$, it survives asymptotically.
\end{frame}

\begin{frame}{Permutation simulation: $L=100$, $k=3$, $m=50$}
\centering
\maybefig[width=0.93\textwidth,height=0.67\textheight,keepaspectratio]{plots/permutation.png}
\end{frame}

\begin{frame}{The two corrections in the permutation simulation}
For $L=100$, $k=3$, exactly 50 successes and 50 failures:
\[
  \bar r \approx 0.4642.
\]

\vspace{0.2cm}
First add the denominator-covariance correction:
\[
  \bar r+\frac{\widehat{\Cov}(r,d)}{\bar d}
  \approx
  0.4642+0.0224
  =0.4866.
\]
This matches the pooled rate:
\[
  \frac{\bar n}{\bar d}\approx 0.4866.
\]

\vspace{0.2cm}
Then add the finite-population correction:
\[
  \frac{k(1-p)}{L-k}
  =
  \frac{3(1/2)}{97}
  \approx 0.0155.
\]
So
\[
  0.4866+0.0155=0.5021\approx 1/2.
\]
\end{frame}

\begin{frame}{A useful comparison: the $SSF$ pattern}
To see that the covariance effect is not just a generic ratio artifact, consider a different pattern:
\[
  SSF \quad \text{followed by } S.
\]
Define
\[
  D=\#\{SSF\text{ patterns with a following trial}\},
  \qquad
  R=\frac{\#\{SSFS\}}{D}.
\]

\vspace{0.2cm}
\centering
\maybefig[width=0.72\textwidth,height=0.43\textheight,keepaspectratio]{plots/ssf.png}
\end{frame}

\begin{frame}{Why the $SSF$ covariance is much smaller}
\small
For $SSS$, a following success reinforces the conditioning event:
\[
  SSS\mapsto SSSS.
\]
It increases the numerator and creates overlapping $SSS$ opportunities.

\vspace{0.1cm}
For $SSF$, the conditioning pattern ends the success run:
\[
  SSF\mapsto SSFS.
\]
The following success increases the numerator, but it does not create the same overlap structure.

\vspace{0.1cm}
In the simulation:
\[
  \widehat{\Cov}(R,D)\approx 0.063,
  \qquad
  \widehat{\Corr}(R,D)\approx 0.199.
\]

\begin{alertblock}{Lesson}
Equal weighting is most dangerous when the conditioning event is reinforced by the outcome being measured.
\end{alertblock}
\end{frame}

\begin{frame}{Main takeaways}
\begin{enumerate}
  \item The bias is not mysterious: it is a denominator--rate covariance.
  \[
    \E[R]=p-\frac{\Cov(R,D)}{\E[D]}.
  \]

  \item Pooling opportunities is the right way to estimate $P(S\mid S^k)$.
  Equal game weights target the mean of ratios, not the conditional probability.

  \item In the fixed-$m$ permutation model, the bias has two pieces:
  \[
    \frac{k(1-p)}{L-k}
    +
    \frac{\Cov(r,d)}{\E[d]}.
  \]

  \item The overlap structure determines the magnitude of the effect.
\end{enumerate}

\vspace{0.2cm}
\begin{block}{Bottom line}
A perfectly memoryless shooter can look colder after a make if we average finite-game conditional percentages in the wrong way.
\end{block}
\end{frame}


\begin{frame}
\frametitle{Problem III: The Paradox of the Blown Lead}
\begin{columns}[T,totalwidth=\textwidth]
\begin{column}{0.50\textwidth}
\begin{block}{The fan interpretation}
If a team reaches win probability $0.78$ and loses, the instinctive reaction is:
\[
  \text{``That should happen only }22\%\text{ of the time.''}
\]
\end{block}

\begin{itemize}
  \item Win-probability graphics invite this interpretation.
  \item But after the game we have selected the team that \emph{lost}.
  \item That changes the reference distribution.
\end{itemize}
\end{column}
\begin{column}{0.48\textwidth}
\centering
\maybefig[width=0.88\linewidth,height=0.55\textheight,keepaspectratio]{plots/blown_espn_superbowl.png}
\end{column}
\end{columns}
\end{frame}



\begin{frame}
\frametitle{The Conditioning Error}
In a two-team game with no ties, the winner eventually reaches win probability $1$.

\vspace{2mm}
For the losing team, define
\[
  M = \max_t WP_{\text{loser}}(t).
\]

\begin{block}{The right comparison}
A losing team that reached $\alpha$ was not a generic team at win probability $\alpha$.
It was drawn from the conditional distribution
\[
  M \mid \{\text{team lost}\}.
\]
\end{block}

\vspace{2mm}
\centering
\Large
\alert{A ``blown lead'' is an extreme of a path, observed after conditioning on loss.}
\end{frame}



\begin{frame}
\frametitle{A Simple Model Where Win Probability is Known}
Assume two evenly matched teams, $A$ and $B$.
\begin{itemize}
  \item Each team has $N$ possessions.
  \item Each possession ends in a score or a turnover.
  \item Each team scores on a possession with probability $1/2$.
  \item If the game is tied after $N$ possessions, overtime is a fair coin flip.
\end{itemize}

Let $d_t=\text{score}_A-\text{score}_B$ and $n_t=N-t$ be the number of remaining possessions. Then
\[
WP_A(t)=\Prb(A\text{ outscores }B)+\frac12\Prb(A\text{ ties }B).
\]

\begin{block}{Why this model is useful}
It gives a controlled experiment where the win-probability model is exactly calibrated by construction.
\end{block}
\end{frame}



\begin{frame}
\frametitle{Simulation Procedure}
With the closed-form win probability in hand:
\begin{enumerate}
  \item Simulate 10,000 games between evenly matched teams.
  \item Compute the true win probability of both teams after each possession.
  \item Identify the eventual loser.
  \item Record the maximum win probability attained by that loser.
\end{enumerate}

\vspace{2mm}
\centering
\maybefig[width=0.78\linewidth,height=0.47\textheight,keepaspectratio]{plots/high-wp-game.png}

\vspace{1mm}
\small One simulated loser reached $98.5\%$ before eventually losing.
\end{frame}



\begin{frame}
\frametitle{The Distribution of the Loser's Maximum Win Probability}
\begin{columns}[T,totalwidth=\textwidth]
\begin{column}{0.56\textwidth}
\centering
\maybefig[width=\linewidth,height=0.54\textheight,keepaspectratio]{plots/blown_sim_distribution.png}
\end{column}
\begin{column}{0.40\textwidth}
\begin{itemize}
  \item The distribution ranges from $0.5$ to $1.0$ for evenly matched teams.
  \item It is decreasing, but not concentrated near $0.5$.
  \item In the simulation, the mean is about $0.690$ and the standard deviation is about $0.139$.
\end{itemize}
\end{column}
\end{columns}

\vspace{1mm}
\begin{block}{Key object}
The object is not $WP(t)$ at a fixed time. It is the maximum of a stochastic process, conditional on eventual loss.
\end{block}
\end{frame}



\begin{frame}
\frametitle{The Threshold Plot: How Often Does the Loser Become a Favorite?}
\begin{columns}[T,totalwidth=\textwidth]
\begin{column}{0.56\textwidth}
\centering
\maybefig[width=\linewidth,height=0.55\textheight,keepaspectratio]{plots/blown_sim_threshold.png}
\end{column}
\begin{column}{0.40\textwidth}
For each threshold $\alpha$, plot
\[
  \Prb(M\ge \alpha).
\]

\begin{itemize}
  \item In half of simulated games, the losing team reached at least $66\%$.
  \item Up to about $\alpha=0.93$, the curve is above the naive line $1-\alpha$.
\end{itemize}
\end{column}
\end{columns}

\vspace{1mm}
\centering
\alert{A losing team being a 2:1 favorite at some point is routine, not shocking.}
\end{frame}



\begin{frame}
\frametitle{Updated NFL Data Look Similar}
\begin{columns}[T,totalwidth=\textwidth]
\begin{column}{0.49\textwidth}
\centering
\maybefig[width=\linewidth,height=0.42\textheight,keepaspectratio]{plots/blown_obs_distribution.png}

\small ESPN NFL regular-season games, $2018$--$2024$, in the $p_0=0.50$ bin: mean $0.676$, SD $0.129$.
\end{column}
\begin{column}{0.49\textwidth}
\centering
\maybefig[width=\linewidth,height=0.42\textheight,keepaspectratio]{plots/blown_obs_threshold.png}

\small In half of observed games, the loser reached at least $67\%$.
\end{column}
\end{columns}

\vspace{2mm}
\begin{block}{Empirical message}
The simple simulation and updated ESPN NFL data tell the same story: high interim win probabilities are common among eventual losers.
\end{block}
\end{frame}



\begin{frame}
\frametitle{Blown-Lead Takeaway}
\begin{itemize}
  \item A win probability model can be calibrated and still produce many dramatic-looking losses.
  \item The correct reference class is not ``teams currently at $\alpha$.''
  \item It is ``teams that ultimately lost, summarized by their maximum win probability.''
\end{itemize}

\vspace{2mm}
\begin{block}{Super Bowl LVII example}
The Eagles reached about $78.4\%$ and lost. The answer to ``how unusual is that?'' is not $21.6\%$.
In the conditional distribution of losing-team maxima, it is about $1/0.784 - 1 \approx 27.6\%$.
\end{block}

\vspace{2mm}
\centering
\Large
\alert{The model need not be wrong; the post-game conditioning changes the question.}
\end{frame}



\begin{frame}
\frametitle{Closing: Three Problems, One Statistical Discipline}
\begin{center}
\small
\begin{tabular}{>{\bfseries}p{0.23\textwidth}p{0.64\textwidth}}
Loser's Curse & Top picks look cursed under expected surplus, but may be rational under a right-tail utility.\\[2mm]
Hot Hand & A memoryless shooter can look cold after makes if we average finite-game conditional percentages in the wrong way.\\[2mm]
Blown Leads & A large maximum win probability is not rare among teams that eventually lose.\\
\end{tabular}
\end{center}

\vspace{4mm}
\begin{block}{Sports analytics at its best}
It turns familiar arguments about games into precise questions about estimands, distributions, reference classes, decisions, and utility functions.
\end{block}
\end{frame}



\begin{frame}
\frametitle{References}
\footnotesize
\begin{itemize}
  \item Cade Massey and Richard H. Thaler. \emph{The Loser's Curse: Decision Making and Market Efficiency in the National Football League Draft}. Management Science, 59(7):1479--1495, 2013.
  \item Ryan S. Brill and Abraham J. Wyner. \emph{The Loser's Curse and the Critical Role of the Utility Function}. arXiv:2411.10400, revised 2025.
  \item Thomas Gilovich, Robert Vallone, and Amos Tversky. \emph{The hot hand in basketball: On the misperception of random sequences}. Cognitive Psychology, 17(3):295--314, 1985.
  \item Joshua B. Miller and Adam Sanjurjo. \emph{Surprised by the hot hand fallacy? A truth in the law of small numbers}. Econometrica, 86(6):2019--2047, 2018.
  \item David M. Ritzwoller and Joseph P. Romano. \emph{Uncertainty in the hot hand fallacy: Detecting streaky alternatives to random Bernoulli sequences}. Review of Economic Studies, 89(2):976--1007, 2022.
  \item Joshua B. Miller and Adam Sanjurjo. \emph{A cold shower for the hot hand fallacy: Robust evidence from controlled settings}. Review of Economics and Statistics, 106(6):1607--1619, 2024.
  \item Jonathan Pipping and Abraham Wyner. \emph{A Paradox of Blown Leads: Rethinking Win Probability in Football}. Working paper, 2026.
  \item Carl S. and Ben Baldwin. \emph{nflfastR: Functions to Efficiently Access NFL Play by Play Data}. R package documentation, 2025.
\end{itemize}
\end{frame}


\end{document}
